\documentclass[11pt,a4paper]{article}

% ============================================================
%  QIMMA -- COURS : Dénombrement
%  2ème Bac Sciences Mathématiques (A et B)
%  Standard exhaustif : définitions formelles + démonstrations
%  Basé sur templates/qimma-template.tex (charte Qimma)
% ============================================================

\usepackage[french]{babel}
\usepackage[utf8]{inputenc}
\usepackage[T1]{fontenc}
\usepackage{lmodern}
\usepackage[a4paper,margin=2cm,top=2.3cm,bottom=2.6cm]{geometry}
\usepackage{amsmath,amssymb}
\usepackage{xcolor}
\usepackage{graphicx}
\usepackage{tikz}
\usetikzlibrary{shadows,calc}
\usepackage[most]{tcolorbox}
\usepackage{titlesec}
\usepackage{fancyhdr}
\usepackage{enumitem}
\usepackage{pifont}
\usepackage{array}
\usepackage{colortbl}
\usepackage{hyperref}

% ------------------- CHARTE QIMMA -------------------
\definecolor{qteal}{HTML}{0d9488}
\definecolor{qtealdark}{HTML}{134e4a}
\definecolor{qamber}{HTML}{fbbf24}
\definecolor{qambertext}{HTML}{92400e}
\definecolor{ink}{HTML}{1f2937}
\definecolor{inkgray}{HTML}{6b7280}
\definecolor{tealpale}{HTML}{e6f5f3}
\colorlet{colDef}{qteal}
\definecolor{colProp}{HTML}{0f766e}
\definecolor{colEx}{HTML}{b45309}
\definecolor{colRem}{HTML}{9d174d}
\color{ink}

\graphicspath{{../../../../assets/}}
\newcommand{\logofile}{../../../../assets/qimma-logo.pdf}

% ------------------- METADONNEES -------------------
\newcommand{\matiere}{Mathématiques}
\newcommand{\niveau}{2\up{ème} Bac -- Sciences Mathématiques}
\newcommand{\chapitretitre}{Dénombrement}

% ------------------- EN-TETE / PIED DE PAGE -------------------
\pagestyle{fancy}
\fancyhf{}
\renewcommand{\headrulewidth}{0pt}
\renewcommand{\footrulewidth}{0.5pt}
\fancyfoot[L]{\raisebox{-0.15cm}{\includegraphics[height=0.35cm]{\logofile}}~\small\sffamily\color{inkgray}Qimma}
\fancyfoot[C]{\small\sffamily\color{inkgray}\matiere\ -- \niveau}
\fancyfoot[R]{\small\sffamily\color{qtealdark}\thepage}

% ------------------- TITRES DE SECTION -------------------
\titleformat{\section}
  {\normalfont\sffamily\Large\bfseries\color{qtealdark}}
  {\tikz[baseline=-3.2pt]{\node[circle,fill=qteal,text=white,inner sep=0pt,minimum size=8mm]{\sffamily\bfseries\thesection};}}
  {10pt}{}
  [\vspace{-2mm}{\color{qamber}\titlerule[1.6pt]}\vspace{2mm}]
\titlespacing*{\section}{0pt}{16pt}{10pt}

\titleformat{\subsection}
  {\normalfont\sffamily\large\bfseries\color{qtealdark}}
  {\color{qteal}\ding{212}}{6pt}{}
\titlespacing*{\subsection}{0pt}{12pt}{6pt}

% ------------------- BOITES PEDAGOGIQUES -------------------
\newtcolorbox{fiche}[3][]{
  enhanced, breakable,
  colback=white, colframe=#2,
  boxrule=1pt, arc=2mm,
  top=7mm, left=4mm, right=4mm, bottom=3mm,
  drop shadow={black!25, shadow xshift=1mm, shadow yshift=-1mm},
  attach boxed title to top left={xshift=4mm, yshift=-3mm, yshifttext=-1mm},
  boxed title style={colback=#2, arc=1.5mm, boxrule=0pt, left=2mm, right=2mm},
  coltitle=white, fonttitle=\sffamily\bfseries,
  title={#3}, #1
}
\newenvironment{definition}[1][Définition]
  {\begin{fiche}{colDef}{\ding{110}~~#1}}{\end{fiche}}
\newenvironment{propriete}[1][Propriété]
  {\begin{fiche}{colProp}{\ding{72}~~#1}}{\end{fiche}}
\newenvironment{exemple}[1][Exemple]
  {\begin{fiche}{colEx}{\ding{217}~~#1}}{\end{fiche}}
\newenvironment{remarque}[1][Remarque]
  {\begin{fiche}{colRem}{\ding{43}~~#1}}{\end{fiche}}

% Démonstration (hors boîte, style discret)
\newenvironment{demo}
  {\par\smallskip\noindent{\sffamily\bfseries\color{qtealdark}Démonstration.}\ \itshape}
  {\ \normalfont\hfill$\blacksquare$\par\medskip}

\hypersetup{colorlinks=true, linkcolor=qtealdark, urlcolor=qtealdark}

% ============================================================
\begin{document}

% ------------------- BANNIERE DE TITRE -------------------
\begin{tcolorbox}[
  enhanced, boxrule=0pt, arc=4mm,
  interior style={left color=qtealdark, right color=qteal},
  top=8mm, bottom=8mm, left=10mm, right=32mm,
  drop shadow={black!35, shadow xshift=1.5mm, shadow yshift=-1.5mm},
  before skip=0pt, after skip=9mm,
  overlay={
    \node[anchor=north west] at ([xshift=6mm,yshift=-6mm]frame.north west)
      {\includegraphics[height=1.25cm]{\logofile}};
    \node[anchor=north east, fill=qamber, text=qambertext, rounded corners=2pt,
          inner sep=4pt, font=\sffamily\bfseries\small] at ([xshift=-6mm,yshift=-6mm]frame.north east) {COURS};
  }
]
\hspace{1.6cm}\centering
{\sffamily\bfseries\color{white}\fontsize{23}{27}\selectfont \chapitretitre}\\[3mm]
{\sffamily\color{white!90}\large \matiere\ \textbullet\ \niveau}
\end{tcolorbox}

% ------------------- OBJECTIFS -------------------
\begin{tcolorbox}[
  enhanced, colback=tealpale, colframe=qteal, boxrule=0.8pt, arc=2mm,
  left=5mm, right=5mm, top=3mm, bottom=3mm,
  title={\sffamily\bfseries\color{qtealdark}\ding{51}~~Objectifs du chapitre},
  coltitle=qtealdark, colbacktitle=tealpale, boxrule=0.8pt,
  attach title to upper={\par\vspace{1mm}}
]
\begin{itemize}[leftmargin=6mm, label=\ding{212}, itemsep=1mm]
  \item Maîtriser le principe multiplicatif et l'utiliser pour dénombrer des situations avec ou sans répétition.
  \item Distinguer et démontrer les formules de $p$-listes, arrangements, permutations et combinaisons.
  \item Calculer avec les factorielles et les coefficients binomiaux, démontrer leurs propriétés (symétrie, formule de Pascal).
  \item Utiliser la formule du binôme de Newton et l'exploiter dans des calculs et démonstrations.
  \item Reconnaître, dans un énoncé, la méthode de dénombrement appropriée (ordre, répétition), avec tous les cas types d'examen.
\end{itemize}
\end{tcolorbox}

\vspace{4mm}

% ------------------- SECTION 1 -------------------
\section{Principe fondamental du dénombrement}

\begin{propriete}[Principe multiplicatif]
Si une situation se décompose en $k$ étapes successives indépendantes, offrant respectivement $n_1,n_2,\ldots,n_k$ possibilités, alors le nombre total de résultats possibles est
\[
  n_1\times n_2\times\cdots\times n_k.
\]
\end{propriete}

\begin{propriete}[Cardinal d'un produit cartésien]
Si $E_1,\ldots,E_k$ sont des ensembles finis, $\text{card}(E_1\times\cdots\times E_k) = \text{card}(E_1)\times\cdots\times\text{card}(E_k)$.
\end{propriete}

\begin{exemple}[Application directe]
Un menu comporte $3$ entrées, $4$ plats, $2$ desserts. Le nombre de menus complets possibles est $3\times4\times2=24$.
\end{exemple}

% ------------------- SECTION 2 -------------------
\section{$p$-listes (tirages avec ordre et répétition)}

\begin{definition}[$p$-liste]
Soit $E$ un ensemble à $n$ éléments. Une \textbf{$p$-liste} (ou $p$-uplet) d'éléments de $E$ est un élément de $E^p$ : une suite ordonnée de $p$ éléments de $E$, \textbf{avec répétition possible}.
\end{definition}

\begin{propriete}[Nombre de $p$-listes]
Le nombre de $p$-listes d'un ensemble à $n$ éléments est $n^p$.
\end{propriete}

\begin{demo}
Chaque position de la liste (il y en a $p$) offre $n$ choix indépendants (répétition autorisée) : par le principe multiplicatif, $n\times n\times\cdots\times n$ ($p$ fois) $=n^p$.
\end{demo}

\begin{exemple}[Application]
Le nombre de mots de $4$ lettres (avec répétition, sur un alphabet de $26$ lettres) est $26^4=456\,976$. Le nombre de codes à $4$ chiffres (chiffres de $0$ à $9$, répétition autorisée) est $10^4=10\,000$.
\end{exemple}

% ------------------- SECTION 3 -------------------
\section{Arrangements (tirages avec ordre, sans répétition)}

\begin{definition}[Arrangement]
Soit $E$ un ensemble à $n$ éléments, $p\leq n$. Un \textbf{arrangement} de $p$ éléments de $E$ est une $p$-liste d'éléments \textbf{deux à deux distincts} de $E$.
\end{definition}

\begin{propriete}[Nombre d'arrangements]
Le nombre d'arrangements de $p$ éléments parmi $n$, noté $A_n^p$, est
\[
  A_n^p = n(n-1)(n-2)\cdots(n-p+1) = \frac{n!}{(n-p)!}.
\]
\end{propriete}

\begin{demo}
Pour le premier élément de l'arrangement : $n$ choix possibles. Pour le deuxième : $n-1$ choix (l'élément déjà choisi est exclu). \ldots Pour le $p$-ième : $n-p+1$ choix. Par le principe multiplicatif : $A_n^p=n(n-1)\cdots(n-p+1)$. En multipliant et divisant par $(n-p)!$, on retrouve $\dfrac{n!}{(n-p)!}$.
\end{demo}

\begin{definition}[Permutation]
Une \textbf{permutation} d'un ensemble à $n$ éléments est un arrangement des $n$ éléments (tous utilisés), soit $A_n^n$.
\end{definition}

\begin{propriete}[Nombre de permutations]
Le nombre de permutations d'un ensemble à $n$ éléments est $n! = n(n-1)(n-2)\cdots2\times1$ (avec $0!=1$ par convention).
\end{propriete}

\begin{exemple}[Applications directes]
\textbf{a)} Nombre de mots de $4$ lettres \textbf{distinctes} pris parmi $26$ : $A_{26}^4=26\times25\times24\times23=358\,800$.

\smallskip
\textbf{b)} Nombre de façons de classer $5$ coureurs à l'arrivée d'une course (sans ex-aequo) : $5!=120$.

\smallskip
\textbf{c)} Nombre de tiercés dans l'ordre parmi $18$ chevaux : $A_{18}^3=18\times17\times16=4\,896$.
\end{exemple}

% ------------------- SECTION 4 -------------------
\section{Combinaisons (tirages sans ordre, sans répétition)}

\begin{definition}[Combinaison]
Soit $E$ un ensemble à $n$ éléments, $p\leq n$. Une \textbf{combinaison} de $p$ éléments de $E$ est une \textbf{partie} (sous-ensemble) de $E$ à $p$ éléments -- l'\textbf{ordre n'intervient pas}.
\end{definition}

\begin{propriete}[Nombre de combinaisons]
Le nombre de combinaisons de $p$ éléments parmi $n$, noté $\dbinom np$, est
\[
  \binom np = \frac{A_n^p}{p!} = \frac{n!}{p!(n-p)!}.
\]
\end{propriete}

\begin{demo}
Chaque combinaison de $p$ éléments donne naissance à $p!$ arrangements distincts (en ordonnant ses éléments de toutes les façons possibles). Donc $A_n^p=\dbinom np\times p!$, d'où $\dbinom np=\dfrac{A_n^p}{p!}=\dfrac{n!}{p!(n-p)!}$.
\end{demo}

\begin{propriete}[Propriétés des coefficients binomiaux]
\[
  \binom n0=\binom nn=1, \qquad \binom n1=n, \qquad \binom np=\binom n{n-p} \quad\text{(symétrie)},
\]
\[
  \binom np+\binom n{p+1}=\binom{n+1}{p+1} \quad\text{(formule de Pascal)}.
\]
\end{propriete}

\begin{demo}[Symétrie]
Choisir $p$ éléments à \textbf{garder} parmi $n$ revient exactement à choisir les $n-p$ éléments à \textbf{écarter} : il y a donc autant de façons de faire l'un que l'autre, soit $\dbinom np=\dbinom n{n-p}$.
\end{demo}

\begin{demo}[Formule de Pascal]
Fixons un élément $a\in E$ ($\text{card}(E)=n+1$). Une partie à $p+1$ éléments de $E$ soit \textbf{contient} $a$ (il reste $p$ éléments à choisir parmi les $n$ restants : $\dbinom np$ façons), soit \textbf{ne contient pas} $a$ (il faut choisir les $p+1$ éléments parmi les $n$ restants : $\dbinom n{p+1}$ façons). Ces deux cas sont disjoints et couvrent tous les cas, d'où $\dbinom{n+1}{p+1}=\dbinom np+\dbinom n{p+1}$.
\end{demo}

\begin{exemple}[Applications directes]
\textbf{a)} Nombre de comités de $3$ personnes parmi $10$ : $\dbinom{10}3=\dfrac{10!}{3!7!}=\dfrac{10\times9\times8}{6}=120$.

\smallskip
\textbf{b)} Nombre de mains de $5$ cartes parmi $32$ : $\dbinom{32}5=201\,376$.

\smallskip
\textbf{c)} Nombre de diagonales d'un polygone à $n$ sommets : $\dbinom n2-n$ (paires de sommets moins les $n$ côtés).
\end{exemple}

% ------------------- SECTION 5 -------------------
\section{Binôme de Newton}

\begin{propriete}[Formule du binôme de Newton]
Pour $a,b\in\mathbb{R}$ (ou $\mathbb{C}$) et $n\in\mathbb{N}$ :
\[
  (a+b)^n = \sum_{k=0}^n\binom nk a^{n-k}b^k.
\]
\end{propriete}

\begin{demo}
Par récurrence sur $n$, en utilisant la formule de Pascal. \textbf{Initialisation} ($n=0$) : $(a+b)^0=1=\dbinom00a^0b^0$. \textbf{Hérédité} : supposons la formule vraie au rang $n$. Alors
\[
  (a+b)^{n+1} = (a+b)\sum_{k=0}^n\binom nk a^{n-k}b^k = \sum_{k=0}^n\binom nk a^{n-k+1}b^k + \sum_{k=0}^n\binom nk a^{n-k}b^{k+1}.
\]
En réindexant la deuxième somme ($k\to k-1$) et en regroupant les termes de même puissance, chaque coefficient devient $\dbinom nk+\dbinom n{k-1}=\dbinom{n+1}k$ (formule de Pascal), ce qui redonne exactement la formule au rang $n+1$.
\end{demo}

\begin{exemple}[Applications du binôme]
\textbf{a)} Développons $(x+1)^4$ :
\[
  (x+1)^4 = \binom40x^4+\binom41x^3+\binom42x^2+\binom43x+\binom44 = x^4+4x^3+6x^2+4x+1.
\]
\textbf{b)} Retrouvons $\displaystyle\sum_{k=0}^n\binom nk=2^n$ en posant $a=b=1$ dans la formule : $(1+1)^n=\displaystyle\sum_{k=0}^n\binom nk=2^n$.

\smallskip
\textbf{c)} En posant $a=1,b=-1$ : $\displaystyle\sum_{k=0}^n(-1)^k\binom nk=0^n=0$ pour $n\geq1$ (alternance des signes, somme nulle).
\end{exemple}

% ------------------- SECTION 6 -------------------
\section{Exercices corrigés -- niveau examen national SM}

\begin{exemple}[Exercice 1 -- distinguer les méthodes]
Une urne contient $5$ boules numérotées de $1$ à $5$. On en tire $3$. Calculer le nombre de tirages possibles dans les cas suivants :
\begin{enumerate}[leftmargin=7mm, itemsep=0.3mm]
  \item tirage successif avec remise (l'ordre compte) ;
  \item tirage successif sans remise (l'ordre compte) ;
  \item tirage simultané (l'ordre ne compte pas).
\end{enumerate}

\smallskip
\textbf{Corrigé.} 1) Avec remise, ordre compte : $p$-liste, $5^3=125$. 2) Sans remise, ordre compte : arrangement, $A_5^3=5\times4\times3=60$. 3) Sans remise, ordre ne compte pas : combinaison, $\dbinom53=10$.
\end{exemple}

\begin{exemple}[Exercice 2 -- anagramme avec contrainte de position]
Combien d'anagrammes du mot MATHS (5 lettres distinctes) commencent par une voyelle ?

\smallskip
\textbf{Corrigé.} MATHS comporte une seule voyelle : A. Fixons la première lettre = A (1 choix). Les $4$ lettres restantes (M,T,H,S) s'arrangent librement sur les $4$ positions suivantes : $4!=24$ façons.

\smallskip
\textbf{Conclusion} : $1\times24=24$ anagrammes commençant par une voyelle.
\end{exemple}

\begin{exemple}[Exercice 3 -- comité avec contrainte de composition]
Un groupe de $12$ personnes comprend $7$ hommes et $5$ femmes. On forme un comité de $4$ personnes. Combien de comités contiennent exactement $2$ hommes et $2$ femmes ?

\smallskip
\textbf{Corrigé.} Choisir $2$ hommes parmi $7$ : $\dbinom72=21$. Choisir $2$ femmes parmi $5$ : $\dbinom52=10$. Par le principe multiplicatif (choix indépendants) :
\[
  21\times10 = 210 \text{ comités}.
\]
\end{exemple}

\begin{exemple}[Exercice 4 -- démonstration avec les coefficients binomiaux]
Démontrer que $\displaystyle\sum_{k=0}^nk\binom nk = n\times2^{n-1}$ pour $n\geq1$.

\smallskip
\textbf{Corrigé.} On part de l'identité $k\dbinom nk=n\dbinom{n-1}{k-1}$ pour $k\geq1$ (vérification : $k\times\dfrac{n!}{k!(n-k)!}=\dfrac{n!}{(k-1)!(n-k)!}=n\times\dfrac{(n-1)!}{(k-1)!(n-k)!}=n\dbinom{n-1}{k-1}$).

\smallskip
Donc
\[
  \sum_{k=0}^nk\binom nk = \sum_{k=1}^nn\binom{n-1}{k-1} = n\sum_{j=0}^{n-1}\binom{n-1}j = n\times2^{n-1}
\]
(en posant $j=k-1$ et en utilisant $\sum\dbinom{n-1}j=2^{n-1}$).
\end{exemple}

\begin{exemple}[Exercice 5 -- au moins / au plus (complémentaire)]
Une urne contient $6$ boules rouges et $4$ boules vertes. On tire simultanément $5$ boules. Combien de tirages contiennent \textbf{au moins une} boule verte ?

\smallskip
\textbf{Corrigé.} Méthode du complémentaire : nombre total de tirages $=\dbinom{10}5=252$. Tirages sans aucune boule verte (5 boules rouges parmi 6) : $\dbinom65=6$.

\smallskip
\textbf{Conclusion} : tirages avec au moins une verte $=252-6=246$.
\end{exemple}

% ------------------- SECTION 7 -------------------
\section{Méthodes et pièges : la boîte à outils de l'examen}

\subsection{Ce qu'on te demande, ce que tu fais}

\begin{center}
\renewcommand{\arraystretch}{1.45}
\sffamily\small
\begin{tabular}{>{\raggedright\arraybackslash}p{4.6cm} >{\centering\arraybackslash}p{2cm} >{\centering\arraybackslash}p{2cm} >{\raggedright\arraybackslash}p{5.6cm}}
\rowcolor{qtealdark} \color{white}\bfseries Situation & \color{white}\bfseries Ordre ? & \color{white}\bfseries Répétition ? & \color{white}\bfseries Formule \\
\rowcolor{tealpale} Code, mot de passe & Oui & Oui & $p$-liste : $n^p$ \\
Tirage successif sans remise & Oui & Non & Arrangement : $A_n^p=\frac{n!}{(n-p)!}$ \\
\rowcolor{tealpale} Classement, rangement complet & Oui & Non & Permutation : $n!$ \\
Comité, tirage simultané & Non & Non & Combinaison : $\binom np=\frac{n!}{p!(n-p)!}$ \\
\end{tabular}
\end{center}

\subsection{Les pièges qui coûtent des points}

\begin{remarque}[Erreurs relevées chaque année par les correcteurs]
\begin{itemize}[leftmargin=6mm, itemsep=0.5mm]
  \item Ne pas se poser la bonne première question : « l'ordre compte-t-il ? », « y a-t-il répétition ? ». Ces deux réponses déterminent \textbf{entièrement} la formule à utiliser.
  \item Confondre $A_n^p$ (arrangement) et $\dbinom np$ (combinaison) : l'arrangement compte $p!$ fois plus (un facteur d'ordre).
  \item Oublier la méthode du \textbf{complémentaire} pour « au moins un » : souvent bien plus simple que de sommer tous les cas favorables.
  \item Dans un problème avec plusieurs catégories (hommes/femmes, rouge/vert), oublier que les choix dans chaque catégorie sont \textbf{indépendants} : on multiplie les combinaisons.
  \item Utiliser une formule d'arrangement/permutation quand il y a \textbf{répétition} possible (auquel cas c'est une $p$-liste, formule $n^p$).
  \item Oublier $0!=1$ dans un calcul de coefficient binomial ou de permutation.
\end{itemize}
\end{remarque}

% ------------------- RESUME -------------------
\vspace{4mm}
\begin{tcolorbox}[
  enhanced, boxrule=0pt, arc=2mm,
  interior style={left color=qtealdark, right color=qteal},
  left=6mm, right=6mm, top=4mm, bottom=4mm,
  fontupper=\color{white}\sffamily,
  title={\sffamily\bfseries\color{white}\ding{72}~~À retenir},
  colbacktitle=qtealdark, coltitle=white, boxrule=0pt
]
\begin{itemize}[leftmargin=6mm, label={\color{qamber}\ding{212}}, itemsep=1mm]
  \item Principe multiplicatif : étapes indépendantes $\Rightarrow$ on multiplie les nombres de choix.
  \item $p$-liste (ordre, répétition) : $n^p$. Arrangement (ordre, sans répétition) : $A_n^p=\frac{n!}{(n-p)!}$. Permutation : $n!$. Combinaison (sans ordre) : $\binom np=\frac{n!}{p!(n-p)!}$.
  \item Symétrie $\binom np=\binom n{n-p}$ ; Pascal $\binom np+\binom n{p+1}=\binom{n+1}{p+1}$.
  \item Binôme de Newton : $(a+b)^n=\sum_{k=0}^n\binom nk a^{n-k}b^k$ ; cas particulier $\sum\binom nk=2^n$.
  \item Méthode du complémentaire pour « au moins un ».
\end{itemize}
\end{tcolorbox}

\end{document}
